\section{Evidence and a research agenda}
\label{sec:research-agenda}

The available evidence supports investigation of the hypothesis, but does not yet test its comparative claim. Murmur's coordination prototype supplies an existence example of published coordination and thread-based context reconstruction. Its commercial gateway remains a specification. The accompanying executable design checks examine the proposed transaction and disclosure rules; those checks isolate necessary controls rather than an advantage of the shared record layer.

The two-deal finite model explores atomic acceptance, seller confirmation or closure, delayed terminal delivery and resource constraints. Its three reference configurations contain 448, 448 and 576 reachable states, respectively, without a violation of the stated invariants. Removing budget or inventory admission permits overcommitment. Releasing funds at timeout permits an already confirmed obligation to lose its backing. These results reproduce from \texttt{protocol\_model.py}; they concern the specified finite abstraction, not an implementation of the gateway, network or cryptography. In particular, no comparison changes the coordination representation while holding enforcement fixed.

The deterministic fixtures provide a similarly limited disclosure check. They reject 13 altered mandate cases and construct identical permitted payloads across 18 work orders and 48 private-context variants per order. The 864 payload checks and 3,456 rejected candidate mutations reproduce from \texttt{boundary\_checks.py}. Private planning context is not an input to the tested release function. The result verifies a fixed interface, not resistance to prompt injection, adaptive information leakage or unauthorised recipient selection. Counts of constructed cases are not estimates of a security failure rate.

\subsection{A comparison that could reject the hypothesis}

The decisive experiment would compare Murmur with signed point-to-point coordination using the same agents, commercial profile, identity trust, gateways, ledgers, release policies and retained evidence. A directory and history retrieval must be available in both designs; omitting them from the comparator would confound representation with missing functionality. Workloads should include stable bilateral negotiation, late participant arrival and delegation across independently operated organisations.

For each workflow, an independent test oracle would enumerate the records required to establish its decision. Reconstruction coverage is the fraction of those required records that an authorised participant can retrieve and verify. Outcome recovery time runs from service resumption or the start of an outcome query until the participant establishes the correct durable terminal state. The experiment should also record unresolved outcomes, duplicate external effects, storage and network cost, and the work needed to integrate a new participant. Thresholds for acceptable cost and meaningful benefit should be fixed before the comparison. No such experiment is reported here.

Failures should be injected at the interfaces that the proposal depends on: after local reservation but before publication, after publication but before acknowledgement, and after seller confirmation but before buyer observation. Duplicate delivery, closure arriving before acceptance, delayed indexing and unavailable historical records require separate cases. A correct recovery must preserve resource backing even when the evidence needed for progress is temporarily unavailable. Reporting successful agent answers alone would miss this distinction.

\subsection{Semantic and disclosure boundaries}

A reproducible commercial test first needs one complete domain profile. The lodging profile must define units, valid quantities, cancellation charges, rounding and a deterministic maximum-exposure calculation. Two independent validators should derive the same obligation from each conformance case and reject unsupported or unbounded liabilities. The illustrative AED~72,000 cap in the worked transaction avoids these ambiguities; it does not resolve them for general trade.

Disclosure evaluation must distinguish permitted business content from private state. Tests should vary private budgets and planning context while keeping approved work orders fixed, then attempt additional fields, value encoding, recipient changes and repeated invocations. Counterparty messages and agent descriptions should be treated as adversarial inputs. Passing a fixed-payload test would support that interface boundary only; adaptive negotiation, timing and cross-deal inference require further policy and analysis.

The interpretation of any positive result must remain conditional. Trusted gateways and serialisable ledgers exclude dishonest resource accounting. Permanent loss of durable state is outside the current recovery model. Amendment accounting and competing replacement revisions require explicit semantics before they enter the supported lifecycle. These are boundaries of the proposal, and each determines which outcome a future experiment can legitimately establish.
